\section{Fourier decay for the actual Jacquet transforms}
\label{jf-section}

We now apply the group operators constructed above to actual
Jacquet integrals.  In particular, the differential operators,
their adjoints, and their decay estimates in this section are
not additional hypotheses.  A separate final paragraph specifies
the further operator identities needed to apply the conclusion
to every recursively generated vector in the source
\cite{alpo-uniform-source}.

\subsection{Smooth Jacquet integrals and their radial adjoints}

Retain the Hermitian spaces, measures and compact representations
of Sections~\ref{rl-section}, \ref{jg-section} and
\ref{jp-section}.  Let \(k\geq1\), \(r\geq0\),
\(G_r=U(k+r,r+1)\), and
\(E=\operatorname{Hom}_{\mathbb C}(W',W)\), with its
Hilbert--Schmidt inner product linear in the first variable.
Let \(J_rF=J_{(1,1)}F\) be the Jacquet integral of
Section~\ref{jp-section}.  Put
\[
 \mathscr S_\varepsilon(G_r,E)=
 \{F\in C^\infty(G_r,E):
       \|F\|_{A,m}<\infty\text{ for all }A,m,
       \ F=0\text{ on }\mathcal G<\varepsilon\},
 \qquad \varepsilon>0.
\]
Here \(\mathcal G\) is the full Gram polynomial
\eqref{rl-Gram}, and \(\|\cdot\|_{A,m}\) is the mixed
left and right derivative seminorm defined in
Section~\ref{jp-section}.  These are the actual coefficient
functions used in the Fourier argument.  Write
\(\mathscr S_+=\bigcup_{\varepsilon>0}\mathscr S_\varepsilon\).

Number the outer coordinates from \(0\) to \(r-1\), and set
\[
 d(c)=\prod_{j=0}^{r-1}c_j,\qquad
 \gamma_r(c)=\prod_{j=0}^{r-1}|c_j|^{2(r-j)+k},\qquad
 d\eta_r(c,m)=\prod_{j=0}^{r-1}|c_j|^{-2}\,dc_j\,dm.
\]
Empty products are one.  The terminal variable \(m\) belongs
to \(G_0=U(k,1)\), restricted to
\(G(m)=|h(mz',z)|^2-TT'>0\).  On this open radial domain put
\begin{equation}\label{jf-radial-transform}
 \Phi_F(c,m)=\gamma_r(c)J_rF(b(c)m).
\end{equation}
The integration formula and Theorem~\ref{jp-pairing-equality}
show that \(\Phi_F\in L^2(d\eta_r,E)\) for every
\(F\in\mathscr S_+\).  More precisely, for fixed compatible
Haar measures there is a positive constant \(C_r\) such that
\begin{equation}\label{jf-radial-pairing}
 \langle\Phi_F,\Phi_H\rangle_{L^2(d\eta_r)}
                          =C_r\beta(F,H).
\end{equation}
Indeed, \(\gamma_r^2\prod|c_j|^{-2}\) is exactly the outer
density in the open-cell Haar formula.  The compact-terminal
integration formula of Section~\ref{jg-section} then gives
\eqref{jf-radial-pairing}; compact equivariance preserves the
Hilbert--Schmidt inner product.  The constant is independent
of \(F,H\), and can be made one by the allowed normalization
of terminal Haar measure.

\begin{lemma}\label{jf-pointwise-smooth}
For \(F\in\mathscr S_+\), the Jacquet integral is absolutely
convergent and smooth at every point of the open cell with
\(\mathcal G>0\).  It may be differentiated there under the
integral sign.  If \(\Theta\) is one of the group operators
of Theorems~\ref{rt-actual-lifts}, \ref{rc-actual-lifts},
or a composition of these with the coefficient of
Theorem~\ref{rl-regular-inverse}, then
\begin{equation}\label{jf-commute}
             \Theta J_rF=J_r(\Theta F),\qquad
             \Theta^\dagger J_rF=J_r(\Theta^\dagger F).
\end{equation}
Both \(\Theta F\) and \(\Theta^\dagger F\) belong to the
same \(\mathscr S_\varepsilon\) as \(F\).
\end{lemma}

\begin{proof}
We give the uniform majorant needed for the pointwise statement.
For one outer pair, the inverse coordinates in
Lemma~\ref{jm-cell-haar} are rational functions of the entries
of the group matrix and of its nonzero lower-left entry.
More explicitly, for
\(g=n_Lw(t)\iota(g')n_R\), let \(J\) be the matrix of
the form on the smaller space.  Then
\[
 t=g_{fe},\quad v_L=g_{We}/t,\quad
 v_R^*J=-g_{fW}/t,\quad
 g'=g_{WW}+v_Lt v_R^*J.
\]
The two central parameters are the imaginary parts of
\(g_{ee}/t\) and \(g_{ff}/t\); their real parts are
determined by isotropy.  Iteration proves that the inverse
coordinates on the full open cell are rational, with
denominators products of the nonzero leading minors and
their conjugates.

Fix a compact subset of the radial domain and allow both
compact factors of the Jacquet integral to vary.  The leading
minors are invariant under the integrating groups and are
bounded away from zero on this compact set.  The preceding
inverse formulas imply
\[
  1+\|n_L\|_{\rm coord}+\|n_R\|_{\rm coord}
              \leq C M(k_L^{-1}n_L^{-1}b(c)m n_Rk_R)^a
\]
for some \(C,a>0\).  Coordinates on the nilpotent groups
have Lebesgue Haar measure.  Every fixed derivative of the
integrand is a finite sum of left or right derivatives of
\(F\), multiplied by polynomials in these nilpotent
coordinates with coefficients bounded on the chosen compact
set.  Choosing the decay exponent in \(\|F\|_{A,m}\)
larger than this polynomial degree and the dimension of the
integration space gives an integrable majorant.  Dominated
convergence proves absolute convergence and every derivative
asserted in the lemma.  It also proves smooth compact
equivariance, either directly or from the smoothness of
finite dimensional continuous unitary representations.

The covariance proved in Sections~\ref{rt-section} and
\ref{rc-section}, and for the spectral coefficient in
Section~\ref{rl-section}, says precisely that these operators
commute with the action \(\rho(p)\).  Moving the operator
inside the now justified integral proves the first identity
in \eqref{jf-commute}.  The action \(\rho(p)\) is unitary
for Haar measure, so its commutation with \(\Theta\)
also implies commutation with \(\Theta^\dagger\).
The polynomial bounds for all derivatives of the coefficients
and their adjoints, on \(\mathcal G\geq\varepsilon\),
prove preservation of every seminorm after increasing its
decay exponent and derivative order.  Differential operators
do not enlarge support.  The resulting functions extend
smoothly by zero across the omitted region, since they vanish
on the open set \(\mathcal G<\varepsilon\).
The same dominated argument proves the second identity.
\end{proof}

\begin{lemma}\label{jf-radial-adjoint}
Suppose a differential operator \(\Theta\) commutes with
the left and right unipotent translations and has radial
operator \(\theta\) after multiplication by \(\gamma_r\).
Then its Haar formal adjoint has radial operator
\(\theta^*\), the formal adjoint for \(d\eta_r\).
This applies to all four group operators used below.
\end{lemma}

\begin{proof}
Use the coordinates \((n_L,c,m,n_R)\) before dividing
by the compact groups.  Haar measure is
\[
 dn_L\,dn_R\,
     \gamma_r(c)^2\prod_j|c_j|^{-2}\,dc_j\,dm.
\]
Take a coordinate chart in \((c,m)\).  In this chart any
finite order differential operator has a unique expansion
in ordered products of a fixed invariant frame on each
nilpotent factor and ordinary coordinate derivatives in
\((c,m)\), with smooth coefficient functions.  This follows
by replacing the ordinary derivatives in the nilpotent
coordinates by the invertible invariant frame and proceeding
from the highest order to the lowest order.  Translation
invariance forces the coefficients in this expansion to be
independent of \(n_L,n_R\).

The nilpotent groups are unimodular.  Each real invariant
vertical vector field consequently has formal adjoint equal
to its negative.  On an equivariant function it acts by a
purely imaginary scalar, because the character is unitary;
evaluation on that character therefore preserves the adjoint
operation.  Horizontal integration by parts uses exactly
the remaining measure
\(\gamma_r^2\prod|c_j|^{-2}dc_j\,dm\).
Thus reduction to the character commutes with formal
adjoints for this measure.  Finally multiplication by the
positive function \(\gamma_r\) is an isometry from that
weighted space to \(L^2(d\eta_r)\), and conjugates both
the operator and its adjoint.  This proves the lemma locally,
hence throughout the open cell.

For the actual operators, both reduced adjoints preserve
the compact equivariance conditions: the group operator
commutes with the unitary compact actions, and the displayed
radial operators and their stars preserve the same conditions.
Their equality on that subspace may therefore be checked
against compactly supported equivariant tests.  Such tests
are obtained by compact averaging, and are dense in the
corresponding local space of sections.  Integration by parts
against them shows that two expressions for the same reduced
operator have the same adjoint on this subspace as well.
\end{proof}

\subsection{A four-operator identity}

Assume now \(r\geq1\).  In the radial coordinates set
\[
 p_j=c_j\partial_{c_j},\quad
 \bar p_j=\bar c_j\partial_{\bar c_j},\quad
 T_j(Y)=\begin{pmatrix}Y+p_j&-c_j\\c_j^{-1}&0\end{pmatrix},
 \quad
 \bar T_j(Y)=
       \begin{pmatrix}Y-\bar p_j&-\bar c_j\\\bar c_j^{-1}&0\end{pmatrix}.
\]
Here the derivatives are the Wirtinger derivatives.
Write
\[
 \binom{A(Y)}{B(Y)}=T_{r-1}(Y)\cdots T_0(Y)\binom10,
 \qquad
 \binom{\bar A(Y)}{\bar B(Y)}=
                  \bar T_{r-1}(Y)\cdots\bar T_0(Y)\binom10.
\]
The bar on \(\bar A,\bar B,\bar T\) is part of this
definition; it does not change the order of differential
operators.  Direct integration by parts for \(|c_j|^{-2}dc_j\)
gives
\[
                    p_j^*=-\bar p_j,\qquad c_j^*=\bar c_j.
\]

The actual terminal matrix of Section~\ref{re-section} is
denoted by \(t(Y)\), with
\(\tau_Y=t_{22}(Y)=Gx(Y)\).  Its finite spectral formula is
\(x(Y)=\sigma\sum_\rho p_\rho(Y)P_\rho\), where
\(\sigma=\pm1\), the \(P_\rho\) are orthogonal projections,
and each \(p_\rho\) is a nonzero real polynomial.  Choose
a real \(Y\) for which all these polynomials are nonzero
at \(Y-1,Y,Y+1\).  Such a choice exists, since there are
only finitely many exceptional real numbers.  Put
\[
 \ell=\tau_Y^{-1}t_{21}(Y),\qquad
 \ell^*=\tau_{Y+1}^{-1}t_{12}(Y+1).
\]
The equality in this notation is the proved shifted adjoint
identity of Section~\ref{re-section}.  These terminal
operators commute with all outer coordinate operators.
Define
\begin{equation}\label{jf-four-radial}
 \begin{split}
 \theta_0&=\bar d\,\bar A(Y+1),\qquad
 \theta_3=d\,A(Y),\\
 \theta_1&=d^{-1}\bigl(\ell A(Y)+B(Y)\bigr),\\
 \theta_2&=\bar d^{-1}
               \bigl(\ell^*\bar A(Y+1)+\bar B(Y+1)\bigr).
 \end{split}
\end{equation}
They are differential operators on the radial domain.

\begin{proposition}\label{jf-four-lifts}
There are actual group differential operators
\(\Theta_0,\Theta_1,\Theta_2,\Theta_3\) with radial
operators \(\theta_0,\theta_1,\theta_2,\theta_3\).
They commute with \(\rho(p)\), and they and their Haar
formal adjoints preserve every \(\mathscr S_\varepsilon\).
For their actual Jacquet transforms, both the operator and
adjoint identities hold pointwise on the radial domain.
\end{proposition}

\begin{proof}
Take \(\Theta_0=\overline{\mathcal A}_{Y+1}\) and
\(\Theta_3=\mathcal A_Y\) from
Theorem~\ref{rt-actual-lifts}.  Theorem~\ref{rc-actual-lifts}
supplies the operators with radial parts
\[
 \bar d\,(t_{21}(Y)A(Y)+\tau_Y B(Y)),\qquad
 d\,(t_{12}(Y+1)\bar A(Y+1)+\tau_{Y+1}\bar B(Y+1)).
\]
Compose on the left with the regular group coefficients
of Theorem~\ref{rl-regular-inverse}, whose radial values
are respectively
\(\mathcal G^{-1}x(Y)^{-1}\) and
\(\mathcal G^{-1}x(Y+1)^{-1}\).
Since \(\mathcal G=d\bar dG\), the results are exactly
\(\theta_1,\theta_2\).  In particular, this construction
does not use \(d^{-1}\) as a coefficient on the full group:
the proved full-group coefficient has only powers of
\(\mathcal G\) in its denominator and is smooth even at
vanishing intermediate minors.  Covariance, growth and
support preservation follow from the three cited theorems
and Leibniz's rule.  Lemmas~\ref{jf-pointwise-smooth} and
\ref{jf-radial-adjoint} give the final assertion.
\end{proof}

For \(s\in\mathbb C\) put
\(f_s(c)=\exp(sc_0-\bar s\bar c_0)\), so that
\(|f_s|=1\).  Multiplication by this function is also
denoted by \(f_s\).

\begin{lemma}\label{jf-four-identity}
On the radial domain one has the exact differential identity
\begin{equation}\label{jf-four-equation}
             \theta_0^*f_s\theta_1
                   -\theta_2^*f_s\theta_3=s f_s.
\end{equation}
\end{lemma}

\begin{proof}
Let \(J=\bigl(\begin{smallmatrix}0&1\\-1&0\end{smallmatrix}\bigr)\).
The relation \([p_j,c_j^{-1}]=-c_j^{-1}\) gives
\[
                  \bar T_j(Y+1)^*JT_j(Y)=J.
\]
The only nonconstant entry to check is the upper-left one:
\[
       (Y+1+p_j)c_j^{-1}-c_j^{-1}(Y+p_j)=0.
\]
For \(j=0\), use \([p_0,f_s]=sc_0f_s\) to obtain
\[
             \bar T_0(Y+1)^*f_sJT_0(Y)=f_sJ+s f_sE_{11}.
\]
For every other \(j\), the wave \(f_s\) commutes with
the entries of \(T_j,\bar T_j\).  Successively cancelling
these factors and taking the upper-left entry yields
\begin{equation}\label{jf-telescoping}
          \bar A(Y+1)^*f_s B(Y)
             -\bar B(Y+1)^*f_s A(Y)=s f_s.
\end{equation}
On substituting \eqref{jf-four-radial}, the left side of
\eqref{jf-four-equation} becomes
\[
 \bar A(Y+1)^*f_s(\ell A(Y)+B(Y))
        -(\bar A(Y+1)^*\ell+\bar B(Y+1)^*)f_sA(Y).
\]
The two terms involving \(\ell\) cancel because \(\ell\)
commutes with \(f_s\).  The remainder is
\eqref{jf-telescoping}.  This calculation retains the
order of every pair of noncommuting operators.
\end{proof}

\subsection{Rapid Fourier decay}

Let \(D_1(g)\) be the first leading minor.  It is a linear
polynomial in the entries of \(g\), invariant under the two
integrating groups, and satisfies \(D_1(b(c)m)=c_0\).
Set
\[
       W_sF(g)=e^{sD_1(g)-\bar s\overline{D_1(g)}}F(g).
\]
For each fixed \(s\), both \(W_s\) and multiplication
by any polynomial in \(D_1,\overline{D_1}\) preserve
\(\mathscr S_\varepsilon\).  This follows by differentiating:
each derivative is the same unit-modulus exponential times
a polynomial in the matrix entries.  Invariance gives
\begin{equation}\label{jf-wave}
                        \Phi_{W_sF}=f_s\Phi_F.
\end{equation}

\begin{theorem}\label{jf-fourier-decay}
For the actual place data and any \(F,H\in\mathscr S_+\),
the function
\[
                       s\longmapsto\beta(W_sF,H)
\]
is a Schwartz function on \(\mathbb C=\mathbb R^2\).
In particular, for every integer \(N\geq0\),
\begin{equation}\label{jf-decay-bound}
        |\beta(W_sF,H)|\leq C_N(F,H)(1+|s|)^{-N}.
\end{equation}
All constants are finite and depend only on the indicated
functions, the fixed place data and the chosen operators.
\end{theorem}

\begin{proof}
Apply \eqref{jf-four-equation} to \(\Phi_F\).  By
Proposition~\ref{jf-four-lifts}, Lemma~\ref{jf-pointwise-smooth}
and \eqref{jf-wave}, its left side is the radial Jacquet
transform of
\[
           \Theta_0^\dagger W_s\Theta_1F
                        -\Theta_2^\dagger W_s\Theta_3F.
\]
Every function here belongs to \(\mathscr S_+\).
Pair with \(\Phi_H\), use \eqref{jf-radial-pairing},
and apply the proved actual adjointness of
Corollary~\ref{jp-poly-adjoint}.  All its hypotheses hold:
the operators commute with \(\rho\), all derivatives of
their coefficients and adjoint coefficients have polynomial
growth on the common positive Gram support, and every input
has derivatives of arbitrarily fast decay.  We obtain
\begin{equation}\label{jf-beta-recursion}
 s\beta(W_sF,H)
       =\beta(W_s\Theta_1F,\Theta_0H)
                    -\beta(W_s\Theta_3F,\Theta_2H).
\end{equation}
No integration by parts on a noncompact radial function has
been assumed in this step; it is supplied by the group
adjointness theorem with its proved integrable majorants.

For the initial uniform bound, use \eqref{jf-wave} and
Cauchy--Schwarz:
\[
 |\beta(W_sF,H)|
 \leq C_r^{-1}\|\Phi_F\|_2\|\Phi_H\|_2.
\]
This bound is independent of \(s\).  Define recursively
\[
 K_0(F,H)=C_r^{-1}\|\Phi_F\|_2\|\Phi_H\|_2,\qquad
 K_{n+1}(F,H)=K_n(\Theta_1F,\Theta_0H)
                           +K_n(\Theta_3F,\Theta_2H).
\]
These numbers are finite, by preservation of \(\mathscr S_+\).
Induction in \eqref{jf-beta-recursion} gives
\( |s|^n|\beta(W_sF,H)|\leq K_n(F,H)\).
Together with the bound for \(n=0\) and
\((1+x)^n\leq2^n(1+x^n)\), this proves
\eqref{jf-decay-bound}.

To include derivatives in \(s\), write \(s=u+iv\) and
\(D_1=x+iy\).  The phase is \(2i(uy+vx)\).
Each real derivative therefore replaces \(F\) by a
polynomial in \(D_1,\overline{D_1}\) times \(F\).
These functions belong to \(\mathscr S_+\).
Their radial transforms are square integrable, so
Cauchy--Schwarz and dominated convergence justify all these
derivatives of the pairing.  Applying the just proved
estimate to each differentiated input proves that every
derivative decays to every order.  This is the Schwartz
assertion.
\end{proof}

\subsection{Compact inputs and smooth lower-rank families}

The following argument provides pointwise regularity of
the marginal densities; a Fourier representative alone
would only give almost everywhere information.
Let \(r\geq1\), and decompose
\(N_r=U_r\rtimes N_{r-1}\), where \(U_r\) is the
one-pair nilpotent group of Lemma~\ref{ju-heisenberg}.
For \(F\in C_c^\infty(G_r,E)\cap\mathscr S_\varepsilon\)
and \(t\ne0\), define a coefficient on \(G_{r-1}\) by
\begin{equation}\label{jf-slice-family}
 \begin{split}
 F_t(g)={}&|t|^{2r+k}
   \int_{U_r\times U_r}
       \psi_z(u_L)\overline{\psi_{z'}(u_R)}\\[-2pt]
       &\hspace{34mm}\cdot
          F(u_L^{-1}w(t)\iota(g)u_R)\,du_L\,du_R.
 \end{split}
\end{equation}
The characters are the restrictions of the original
unipotent characters.  Use recursively compatible Haar
measures; other choices only insert fixed positive constants.

\begin{proposition}\label{jf-compact-family}
The function \((t,g)\mapsto F_t(g)\) is smooth on
\(\mathbb C^\times\times G_{r-1}\).  On every compact
\(K\subset\mathbb C^\times\), all \(F_t\) have support
in one compact subset of \(G_{r-1}\), and there is
\(\varepsilon_K>0\) such that \(F_t=0\) on
\(\mathcal G_{r-1}<\varepsilon_K\) for every \(t\in K\).
All their parameter derivatives have the same two properties.
The exact radial factorization is
\begin{equation}\label{jf-factorization}
                  \Phi_F(t,c',m)=\Phi^{(r-1)}_{F_t}(c',m).
\end{equation}
\end{proposition}

\begin{proof}
Apply the inverse coordinate formulas in the proof of
Lemma~\ref{jf-pointwise-smooth} to
\(u_L^{-1}w(t)\iota(g)u_R\) in \(\operatorname{supp}F\).
When \(t\) remains in \(K\), its modulus has a positive
lower bound.  All inverse coordinates \(u_L,u_R,g\) are
then bounded uniformly.  They lie in closed subsets of
finite dimensional matrix spaces, so they lie in fixed
compact sets.  Thus the integrand is supported in one
compact subset of the integration variables and \(g\),
locally uniformly in \(t\).  Differentiation under this
compact integral proves smoothness and the support assertion.
Derivatives do not enlarge this support; taking a slightly
larger compact neighborhood of \(K\) gives the assertion
also at its boundary.

Leading minors in these coordinates satisfy
\(D_{j+1}^{(r)}=tD_j^{(r-1)}\).  In particular,
\[
 \mathcal G_r(u_L^{-1}w(t)\iota(g)u_R)
                                  =|t|^2\mathcal G_{r-1}(g).
\]
Consequently one may take
\(\varepsilon_K=\varepsilon/\max_{t\in K}|t|^2\).
For derivatives, use a compact neighborhood of \(K\) in
this formula.  A smooth function which vanishes on the
corresponding open set has all derivatives zero there.

The semidirect-product Haar formula and the character
identity factor the full Jacquet integral into the outer
integral in \eqref{jf-slice-family} followed by the smaller
Jacquet integral.  To check the ordering, move the two
\(N_{r-1}\) variables and the compact factors across
\(w(t)\), which commutes with the embedded smaller group.
Conjugate \(u_L,u_R\) by those factors.  Their action on
\(U_r\) preserves Haar measure: the unipotent factors
have determinant one, and the compact factors act unitarily.
It preserves the restricted characters because the original
character is a homomorphism and is invariant under the
compact stabilizers.  No Jacobian or phase remains.
All integrals at a regular cell point have compact support
in their nilpotent coordinates, so Fubini is legitimate.
Finally
\(\gamma_r(t,c')=|t|^{2r+k}\gamma_{r-1}(c')\),
which proves \eqref{jf-factorization} including its scalar
normalization.
\end{proof}

\begin{corollary}\label{jf-smooth-slices}
For compactly supported \(F\) as above, the map
\[
       t\longmapsto\Phi_F(t,\cdot,\cdot)
       \quad\text{from }\mathbb C^\times
                  \text{ to }L^2(d\eta_{r-1},E)
\]
is smooth.  Its derivatives are its pointwise derivatives.
The same holds after any finite word in
\[
       t\partial_t,\quad\bar t\partial_{\bar t},\quad
             t,\quad t^{-1},\quad\bar t,\quad\bar t^{-1},
\]
and after multiplication by any fixed wave
\(e^{st-\bar s\bar t}\).  All corresponding Hilbert norms
are bounded locally uniformly for \(t\ne0\).
\end{corollary}

\begin{proof}
On a compact parameter neighborhood, the family in
Proposition~\ref{jf-compact-family} and all its derivatives
have common compact support and a common positive Gram
bound.  Taylor's formula with its integral remainder,
differentiated a fixed number of times on the smaller group,
shows that the family is smooth for every seminorm
\(\|\cdot\|_{A,m}\).  The Jacquet map into
\(L^2(d\eta_{r-1},E)\) is continuous for
\(\|\cdot\|_{A,4}\), by
Theorem~\ref{jp-pairing-equality} and its bound.
It therefore sends this smooth family to a smooth Hilbert
space valued family, with the asserted derivatives.
Pointwise differentiation is also valid by the compact
integral argument.  All six displayed operators and the
wave act on the parameter by smooth coefficients on
\(\mathbb C^\times\), which proves the remaining claims.
\end{proof}

\begin{theorem}\label{jf-marginal-schwartz}
For \(F,H\in\mathscr S_+\), the integrable marginal density
\begin{equation}\label{jf-density}
 D_{F,H}(t)=|t|^{-2}
     \int\langle\Phi_F(t,c',m),\Phi_H(t,c',m)\rangle
                                    \,d\eta_{r-1}(c',m)
\end{equation}
has a Schwartz representative on \(\mathbb C\).
If \(F,H\) are compactly supported, the integral in
\eqref{jf-density} exists for every \(t\ne0\), is smooth
there, and agrees there with that representative.  It
therefore has a unique smooth extension across \(t=0\),
and that extension and all its derivatives decay faster
than every inverse power at infinity.
\end{theorem}

\begin{proof}
Fubini and Cauchy--Schwarz show that the density exists
almost everywhere and belongs to \(L^1(\mathbb C)\).
Its Fourier transform, using the real bilinear phase
\(2\operatorname{Im}(st)\), is
\[
 \int_{\mathbb C}e^{st-\bar s\bar t}D_{F,H}(t)\,dt
                             =C_r\beta(W_sF,H).
\]
Theorem~\ref{jf-fourier-decay} makes this a Schwartz
function.  Ordinary Fourier inversion gives a Schwartz
representative of \(D_{F,H}\).  To recall why the assertion
includes every weight and derivative, differentiation of
the inverse transform inserts a polynomial in the Fourier
variable, and multiplication by a coordinate is integration
by parts in that variable.  All resulting integrals converge
absolutely for a Schwartz Fourier transform.  Approximate
identities identify the inverse transform with the original
\(L^1\) density almost everywhere.

For compact \(F,H\), Corollary~\ref{jf-smooth-slices}
shows that the inner product in \eqref{jf-density} is a
smooth function of \(t\ne0\), with all derivatives given
by the finite Leibniz sums of Hilbert inner products.
It is therefore continuous and agrees almost everywhere
with the smooth Fourier representative.  Two continuous
functions on an open subset of the plane which agree almost
everywhere agree everywhere: otherwise their difference
would stay nonzero on a ball of positive measure.
This proves the claimed pointwise extension and all its
decay properties.
\end{proof}

\subsection{Terminal support and the remaining application}

For compactly supported \(F\), terminal support is also
locally compact in the outer coordinates.  Indeed, on any
compact subset of \((\mathbb C^\times)^r\), the quantity
\(|d(c)|\) is bounded above and away from zero.  The invariant
minor \(D_{r+1}=d(c)\zeta\) is bounded on
\(\operatorname{supp}F\), and
\(\mathcal G=|d(c)|^2G\geq\varepsilon\) on any contributing
orbit.  Thus \(|\zeta|\) is bounded and \(G\) is bounded
away from zero.  By the compact double-coset description
in Section~\ref{jg-section}, the allowed terminal matrices
lie in a compact subset of \(G_0\).  Every terminal
differential operator and every smooth scalar multiplier
preserves this support.  It follows that all finite terminal
operator words applied to these vectors belong to terminal
\(L^2\), and their actual Haar formal adjoints may be used
by ordinary compactly supported integration by parts.
The assertion includes \(r=0\), when there are no outer
variables and no Fourier-decay statement is required.

The source's generated compact class starts with these
vectors, applies finite differential words in the current
outer coordinate when passing to the next level, and
multiplies by unit-modulus waves.  Repeating
Proposition~\ref{jf-compact-family} proves its actual local
slice regularity and terminal compactness at every fixed
nonzero outer parameter.  This also proves the local
Leibniz formulas for its partial pairings.  More explicitly,
after fixing the first \(i\) nonzero coordinates, inductively
represent the vector as the radial Jacquet transform of a
compactly supported smooth coefficient on \(G_{r-i}\).
The coefficients depend smoothly on those fixed parameters,
with common compact support and positive Gram bound on
each compact parameter set.  A wave in the next coordinate
is multiplication of that coefficient by the corresponding
wave in its first leading minor, which preserves these
properties.  Every entry of the source's matrix \(bT\)
is a finite word in the six displayed coordinate operators.
Applying such an entry and then fixing that coordinate is
handled by differentiating the family
\eqref{jf-slice-family}; Proposition~\ref{jf-compact-family}
and Corollary~\ref{jf-smooth-slices} preserve exactly the
inductive properties.  Finite linear combinations do so
as well.  At \(i=r\), a terminal wave is a smooth multiplier
on the already compact terminal support.  This proves the
claim by induction, including the empty sequence of steps.

For the full residue matrices, one further algebraic step
must be kept distinct: all four entries of the initial
matrix \(Q_0(Y)\) must be identified with their full-group
operators, and that identification must intertwine the
successive one-pair restrictions.  The four operators
\(\Theta_j\) proved above do not by themselves establish
that larger matrix identity.  Once it is established, its
polynomial differential operators preserve compact support
and the positive Gram cutoff, so the present theorems apply
to every finite product of those entries.  Until that
identification is supplied, this section proves the actual
Fourier decay, pointwise smooth densities, generated local
slice regularity and terminal integrability stated above;
it does not assert all of the remaining residue-pairing
identities by assumption.


